\scanpage{51}
% typed verso
% @top right, pagination
\mspage{5}
% @01–04, continuation
algebra of $X$ (in any sense), then one finds the second homotopy \unsure{grou|group} $\pi_2 X$ has rank $n$, $\pi_3 X$ has rank $\ell+m$, the Hurewicz map $\pi_3 X\to H_3 X$ has rank $\ell$, and the symmetric pairing given by Whitehead product
% @05, bracket
\[
 \pi_2 X\handwritten{\otimes}\pi_2 X\mathrel{\overset{[,\!]}{\handwritten{\longrightarrow}}}\pi_3 X
\]
% @06–07, tensor
is described over the rationals by the $i,j$ symmetric tensor $(a_{ij}^k)$.)

% @08, induction
Now suppose by induction that
% @09, model and map
\[
 \Lambda(x_1,\ldots,x_n;y_1,y_2,\ldots;z_1,z_2,\ldots:d)
 \mathrel{\overset{\rho_n}{\handwritten{\longrightarrow}}}\handwritten{\mathscr{E}}
\]
% @10–13, induction hypothesis
has been constructed so that $d$ of any generator is a polynomical in previous generators while $\rho_n$ is an isomorphism on cohomology up through dimension $n$ and injective in dimension $n+1$.

% @14–25, adjoining variables
Then as in step 2 we adjoin variables in dimension $n+1$ to make $\rho$ onto in dimension $n+1$ and injective in dimension $n+2$. In this general case we add variables $\omega_1,\ldots,\omega_d,\omega_{d+1},\ldots,\omega_r$ and define $d$ and $\rho=\rho_{n+1}$ as above. Now $d\omega_1,\ldots,d\omega_d$ are zero and $\rho\omega_1,\ldots,\rho\omega_d$ give generators of the cokernel of $\rho_n$ on $H^{n+1}$, $d\omega_{d+1},\ldots,d\omega_r$ are polynomials in $x_1,x_2,\ldots;y_1,y_2,\ldots;z_1,z_2,\ldots$ which generate the kernel of $\rho_n$ on $H^{n+2}(\Lambda(x_1,x_2,\ldots;y_1,y_2,\ldots;z_1,z_2,\ldots:d))$ and $\rho\omega_{d+1},\ldots,\rho\omega_r$ are elements in $\mathcal{E}$ whose differentials are these same polynomials in
% @26, final display
\[
 \rho x_1,\rho x_2,\ldots;\rho y_1,\rho y_2,\ldots;\rho z_1,\rho z_2,\ldots\ .
\]
% archive sheet number: 50

